% New, unfrozen INPUT FRAGMENT, 18 October 2026.
% Uses the theorem environments, \HH, \ip, \OCT, \DBL and TikZ libraries
% of ideal_ra4_octahedral_decomposition.tex.  It has not been built.
% No existing manuscript or PDF is modified by this file.

\section{An extension excluding square-antiprism facets}
\label{sec:a4extension}

Let \(\mathcal A\) denote the combinatorial square antiprism: it has
eight vertices, two square faces, and eight triangular faces. We first
determine the intrinsic geometry of any ideal right-angled realization.
The proof is elementary and does not invoke a general rigidity theorem.

\begin{lemma}\label{lem:a4rigid}
Every ideal right-angled hyperbolic square antiprism is isometric to
the standard symmetric one. At each cusp its rectangular link has
aspect ratio \(\sqrt2:1\). Exactly one side is quadrilateral, and
that side is long. A triangular face occupies a short side at the
two vertices on one square base, and a long side at its third vertex.
\end{lemma}
\begin{proof}
Use Lorentz coordinates with product
\(-tt'+xx'+yy'+zz'\). Distinct outward unit normals have product
zero or at most \(-1\). To recall the gap argument, for a product
\(g\in(-1,0)\) the two-wall Gram matrix is positive definite
with nonnegative inverse. Projection of a strict timelike interior
point to both walls adds positive wall coefficients and keeps every
other inequality strict by acuteness. This would give an actual
nonright ridge, a contradiction. Let \(b,c\) be the normals of
the two square bases. They cannot be orthogonal: acute orthogonal
projection would then produce an actual ridge. They cannot have
product \(-1\): the future null sum \(b+c\) would satisfy every
other wall inequality by acuteness, producing an actual shared cusp.
The bases have disjoint vertex sets and no shared ridge, so their
planes are ultraparallel. We may write
\[
 b=(0,0,0,1),\qquad c=(\sinh L,0,0,-\cosh L),\qquad L>0.
\]
The four upper cusps, in cyclic order, can be normalized as
\(\xi_i=(1,\cos\theta_i,\sin\theta_i,0)\).
The lower lateral triangle containing \(\xi_i\) has normal
\(\ell_i\) opposite \(b\) in that rectangular cusp link.
Thus \(\ell_i=\mu_i\xi_i-b\) with \(\mu_i>0\).
This triangle also meets the lower square base orthogonally, giving
\[
 0=\ip{\ell_i}{c}=-\mu_i\sinh L+\cosh L,
 \qquad \mu_i=\coth L=:k>1.
\]
Consecutive \(\ell_i,\ell_{i+1}\) meet at a lower cusp as an
opposite pair. Their product is therefore \(-1\), so
\[
 -1=k^2\bigl(\cos(\theta_{i+1}-\theta_i)-1\bigr)+1.
\]
All four positive cyclic gaps are \(\alpha\) or \(2\pi-\alpha\)
for the same \(\alpha\in(0,\pi)\). A gap of the latter kind
would make their sum at least \(2\pi+2\alpha\).
Hence all four equal \(\pi/2\), and \(k=\sqrt2\).
This fixes the regular upper base, the four normals \(\ell_i\),
and the lower base. Each remaining triangle plane is fixed by an
upper base edge and orthogonality to \(b\). All ten supporting
planes are determined, proving rigidity.

For the link calculation put
\[
 h=\sqrt2-1,\qquad r=\sqrt{2h},\qquad r^2=1-h^2.
\]
A model has upper and lower vertices
\begin{align*}
 v_j&=(1,r\cos(j\pi/2),r\sin(j\pi/2),h),\\
 w_j&=(1,r\cos(j\pi/2+\pi/4),r\sin(j\pi/2+\pi/4),-h),
 \qquad j=0,1,2,3.
\end{align*}
The base normals are \(B_\pm=(h/r,0,0,\pm1/r)\).
At \(v_0\), the remaining three normals are
\[
 N_+=(1/r,1,1,h/r),\quad
 N_-=(1/r,1,-1,h/r),\quad
 M=(1/r,\sqrt2,0,-h/r).
\]
They are unit, their adjacent products vanish, and their opposite
products equal \(-1\). The opposite sums are
\[
 N_++N_-=\frac2r v_0,\qquad
 B_++M=\frac{\sqrt2}{r}v_0.
\]
In upper-half-space, the sum of opposite vertical unit wall normals
is their Euclidean separation times one common cusp null vector.
Their coefficients therefore give the two rectangle widths.
The side on \(B_+\) spans the larger width and is long.
It is the only quadrilateral side. Symmetry carries this calculation
to every vertex.

The triangle \(T=(v_0,v_1,w_0)\) has normal \(N_+\).
At \(v_0,v_1\) it is a short rectangle side. At \(w_0\) the
opposite sum is
\[
 B_-+N_+=\frac{\sqrt2}{r}w_0,
\]
whereas the other opposite pair has coefficient \(2/r\).
Here \(T\) spans the larger width and is long, as asserted.
\end{proof}

\begin{figure}[ht]
\centering
\begin{tikzpicture}[scale=.83,line cap=round,line join=round]
 \begin{scope}
  \fill[blue!4] (0,0) rectangle(3.15,2.2);
  \draw[very thick,blue!65!black] (0,2.2)--(3.15,2.2);
  \draw[very thick,red!70!black] (0,0)--(3.15,0);
  \draw[thick] (0,0)--(0,2.2) (3.15,0)--(3.15,2.2);
  \node[above,align=center] at(1.575,2.2)
   {square base\\one long quadrilateral side};
  \node[below,align=center] at(1.575,0)
   {opposite long triangular side\\neighbor in the belt};
  \node[rotate=90,above] at(0,1.1){short triangle};
  \node[rotate=90,below] at(3.15,1.1){short triangle};
  \draw[{Stealth}-{Stealth}] (0,-1.05)--node[below]{\(\sqrt2\)}(3.15,-1.05);
  \draw[{Stealth}-{Stealth}] (3.7,0)--node[right]{\(1\)}(3.7,2.2);
  \node[align=center] at(1.575,-2.0)
   {intrinsic cusp rectangle of \(\mathcal A\)};
 \end{scope}
 \begin{scope}[xshift=6.8cm,yshift=.1cm]
  \coordinate(a) at(0,0); \coordinate(b) at(3.3,0);
  \coordinate(c) at(3.3,1.7); \coordinate(d) at(0,1.7);
  \coordinate(aa) at(.8,.65); \coordinate(bb) at(4.1,.65);
  \coordinate(cc) at(4.1,2.35); \coordinate(dd) at(.8,2.35);
  \fill[blue!13] (a)--(d)--(dd)--(aa)--cycle;
  \fill[blue!13] (b)--(c)--(cc)--(bb)--cycle;
  \draw[thick] (a)--(d)--(dd)--(aa)--cycle;
  \draw[thick] (b)--(c)--(cc)--(bb)--cycle;
  \draw[very thick,red!70!black] (a)--(b) (dd)--(cc);
  \draw[thick,dashed] (d)--(c) (aa)--(bb);
  \node[fill=white,inner sep=1pt] at(.35,1.2){\(\OCT\)};
  \node[fill=white,inner sep=1pt] at(3.75,1.2){\(\OCT\)};
  \node at(1.9,.85){\(\mathcal A\) belt};
  \node[above] at(2.45,2.35){\(\mathcal A\) belt};
  \draw[{Stealth}-{Stealth}] (0,-.6)--node[below]{\(\sqrt2\)}(3.3,-.6);
  \node[align=center] at(2,-1.95)
   {ambient \(\sqrt2:1:1\) box\\
    two opposite long edges are quadrilateral\\
    every end edge is triangular};
 \end{scope}
\end{tikzpicture}
\caption{Cusp-side data and the only noncubic box that can contain an
antiprism cusp. Red edges are quadrilateral-ridge incidences; the other
edges are triangular-ridge incidences. The box has four antiprism
rectangles around its belt and octahedral square ends. The displayed
length ratio is schematic. This is a necessary local box pattern,
not a claimed complete four-polytope.}
\label{fig:a4box}
\end{figure}

\begin{theorem}\label{thm:a4exclude}
Suppose every facet of an entirely ideal finite-volume right-angled
four-polytope is isometric to \(\OCT\) or \(\DBL\), or is
combinatorially a square antiprism. Then no square-antiprism facet
occurs.
\end{theorem}
\begin{proof}
Recall the intrinsic signatures of the first two types. An
octahedral cusp rectangle is square and all four sides are triangular.
A doubled-octahedral rectangle is either square with exactly one
quadrilateral side, or \(2:1\) with both long sides quadrilateral
and both short sides triangular. Lemma~\ref{lem:a4rigid} adds
aspect \(\sqrt2:1\), with one long quadrilateral side.

At an ambient cusp order the box lengths \(x\ge y\ge z\)
and normalize \(z=1\). All pairwise ratios belong to
\(\{1,\sqrt2,2\}\), so the complete list is
\[
 1:1:1,\quad \sqrt2:1:1,\quad\sqrt2:\sqrt2:1,\quad
 2:1:1,\quad2:\sqrt2:1,\quad2:2:1.
\]
The triangle or quadrilateral label on a box edge must agree in
its two incident facet rectangles. In \(2:2:1\), all four
edges on each square end are forced quadrilateral by the doubled
belt faces, exceeding the square-end capacity of one.
In \(\sqrt2:\sqrt2:1\), four antiprism belt faces each
require one long quadrilateral edge on the two square ends;
the total demand four exceeds their total capacity two.
In \(2:\sqrt2:1\), the two \(2:1\) doubled faces force
both long edges of each \(2:\sqrt2\) antiprism face to be
quadrilateral, also impossible.

The remaining \(2:1:1\) box contains four doubled belt facets
and two octahedral ends. The cube contains no antiprism rectangle.
Thus an antiprism cusp must be in a \(\sqrt2:1:1\) box.
Its four belt faces are antiprism rectangles. All their short sides
are triangular, so both square ends have no quadrilateral incidence
and are octahedral. Their long quadrilateral sides form a matching
on the four-face belt. In particular a short side of an antiprism
rectangle meets an octahedral end, and a long side meets an
antiprism belt facet.

Now let \(F\) be an antiprism facet and \(T\) one of its lateral
triangles. Across the relative interior of this actual triangular
ridge there is one other facet \(F'\). At the two vertices of
\(T\) on one square base, Lemma~\ref{lem:a4rigid} makes
\(T\) a short cusp side, so the box forces \(F'\) to be
octahedral. At its third vertex \(T\) is a long cusp side,
so the box forces the same \(F'\) to be an antiprism.
These are distinct facet types, a contradiction.
\end{proof}

\begin{figure}[ht]
\centering
\begin{tikzpicture}[line cap=round,line join=round,scale=.95]
 \coordinate(v0) at(-2,2); \coordinate(v1) at(2,2);
 \coordinate(w0) at(0,-.25);
 \draw[thick] (v0)--node[above]{edge on a square base}(v1)--(w0)--cycle;
 \foreach \p/\lab/\pos in {v0/v_0/above left,v1/v_1/above right,w0/w_0/below}{
  \node[circle,draw,fill=white,inner sep=2.7pt,label={\pos:\(\lab\)}] at(\p){};
 }
 \node at(0,1.2){actual triangular ridge \(T\)};
 \node[draw,rounded corners,align=center,fill=blue!4] (F) at(0,-2.1)
  {one neighboring facet \(F'\) across \(T\)\\
   must have one fixed isometry type};
 \begin{scope}[xshift=-4.4cm,yshift=-.5cm]
  \draw[thick] (0,0) rectangle(1.6,1.05);
  \draw[very thick,red!70!black] (1.6,0)--(1.6,1.05);
  \node[align=center] at(.8,-.65){short at \(v_0\)\\
   \(F'\) is an \(\OCT\) end};
  \draw[-{Stealth},dashed] (1.6,.9)--(v0);
 \end{scope}
 \begin{scope}[xshift=2.8cm,yshift=-.5cm]
  \draw[thick] (0,0) rectangle(1.6,1.05);
  \draw[very thick,red!70!black] (0,0)--(0,1.05);
  \node[align=center] at(.8,-.65){short at \(v_1\)\\
   \(F'\) is an \(\OCT\) end};
  \draw[-{Stealth},dashed] (0,.9)--(v1);
 \end{scope}
 \draw[thick] (-.8,-3.6) rectangle(.8,-2.65);
 \draw[very thick,red!70!black] (-.8,-2.65)--(.8,-2.65);
 \node[align=center] at(0,-4.2){long at \(w_0\)\\
  \(F'\) is an \(\mathcal A\) belt facet};
\end{tikzpicture}
\caption{The global consistency that a formal cusp-box coloring omits.
The highlighted side represents the same triangular ridge at each of
its three ideal vertices. Intrinsic antiprism geometry gives the
short, short, long signature. The box would require its one neighboring
facet to be both octahedral and antiprismatic. The small rectangles
are link schematics; no four-dimensional example is asserted.}
\label{fig:a4triangle}
\end{figure}

\begin{corollary}\label{cor:a4chain}
The octahedral decomposition and chain theorem applies with the facet
hypothesis enlarged to allow combinatorial ideal right-angled square
antiprisms. They do not occur, so the class returns to the original
octahedral and doubled-octahedral spectrum. All its volume and uniqueness
conclusions retain that restricted scope.
\end{corollary}

The exclusion itself uses only the displayed cusp signatures and the
actual neighboring facet across a triangle. Additional facet types
with the same doubled-octahedral cusp signatures would not invalidate
this local exclusion, provided only octahedra have a square cusp with
zero quadrilateral sides. This observation does not extend the global
surface or chain theorem to those types. A larger octahedral gluing
can contain several internal triangles or a new cusp ratio, and the
global argument would need to be proved again. In particular its
path-rigidity step uses the upper bound of two octahedra in a boundary
facet; permitting three removes that contradiction.

\begin{corollary}\label{cor:a4ideal5}
Every entirely ideal finite-volume right-angled five-polytope has
a three-face outside the spectrum \(\{\OCT,\DBL,\mathcal A\}\).
The same holds in dimension six.
\end{corollary}
\begin{proof}
For an ideal right-angled four-polytope \(Q\), put
\(\kappa(Q)=3f_0(Q)-4f_3(Q)+24\).
An \(m\)-cell chain has \(f_0=18m+6,\ f_3=14m+10\), hence
\(\kappa=2-2m\le0\).

We recall the universal five-dimensional identity directly.
In an ideal right-angled five-polytope write \(v,e,r,g,f\)
for the numbers of faces of dimensions zero through four.
Its cusp four-box gives \(e=8v\), twenty-four three-faces per
cusp, and eight four-facets per cusp. An ideal right-angled
three-face has two more polygon faces than vertices.
Since each two-face belongs to three three-faces,
\[
 3r=24v+2g.
\]
Euler's identity \(v-e+r-g+f=2\) then gives
\(g=3v+3f-6\). Every three-face belongs to two four-facets,
so
\[
 \sum_Q\kappa(Q)=3(8v)-4(2g)+24f=48.
\]
If all three-faces belonged to the stated spectrum, each four-facet
would be a chain by Corollary~\ref{cor:a4chain} and would have
nonpositive \(\kappa\), contradicting this sum. An ideal
six-polytope has ideal five-facets, to which the result applies.
\end{proof}

This corollary excludes an explicit three-face spectrum, not a
dimension. A numerical reformulation as a three-face with at least
ten vertices requires the separate complete classification of ideal
right-angled three-polyhedra with at most nine vertices. We neither
prove nor invoke that classification in this appendix.

This appendix is a new written proof. Its symbolic calculations were
not executed by a checker or formalized in Lean. It makes no claim
about all ideal right-angled four-polytopes, universal volume
divisibility, or worldwide priority. The original manuscript and
its compiled PDF remain unchanged.
